% Minimal preamble for Open Logic Project sources (pinned commit 1e960beff9ed7835bf3e3f1335e21af3439cd107).
%
% OLP's own master preamble cannot run on the project host: TeX Live 2025 there lacks
% prooftree.sty, ntheorem.sty, extarrows.sty, mathdots.sty, isomath.sty and subfiles.sty.
% Following the project's own rule ("if the author preamble needs packages the toolchain lacks,
% do not invent macros -- use a minimal working preamble"), this file declares, verbatim in
% spirit and numbering behaviour, exactly the theorem-like environments OLP defines in
% sty/open-logic-book-envs.sty (thm/lem/prop/cor/defn/ex/prob/rem/note/case/axiom/conv),
% plus the shorthand accents OLP uses (`!!a`, `!!b`, `!!{...}`).
%
% Author text is untouched; only the packaging is ours.
\documentclass[10pt,a4paper]{article}

\usepackage{amsmath}
\usepackage{amssymb}
\usepackage{amsthm}
\usepackage{graphicx}
\usepackage{hyperref}
\usepackage{enumitem}
\usepackage[framemethod=default]{mdframed}

% ---- OLP shorthand accents (from sty/open-logic.sty) ----
% NB: \! is taken by LaTeX (negative thin space), so we only define the !!a / !!b families
% and the \olup helper used by them.
\newcommand*{\olup}[1]{\MakeUppercase{#1}}
\def\!!!a#1{\olup{#1}}
\def\!!a#1{\olup{#1}}
\def\!!b#1{\olup{#1}}

% ---- theorem-like environments, same family numbering as OLP ----
\theoremstyle{plain}
\newtheorem{thm}{Theorem}[section]
\newtheorem{lem}[thm]{Lemma}
\newtheorem{prop}[thm]{Proposition}
\newtheorem{cor}[thm]{Corollary}
\newtheorem{axiom}{Axiom}

\theoremstyle{definition}
\newtheorem{defn}[thm]{Definition}
\newtheorem{ex}[thm]{Example}
\newtheorem{prob}{Problem}[section]
\newtheorem{probd}{Problem}
\newtheorem{conv}[thm]{Convention}

\theoremstyle{remark}
\newtheorem{rem}[thm]{Remark}
\newtheorem{note}[thm]{Note}
\newtheorem{case}[thm]{Case}

% ---- environments OLP defines as wrappers ----
\newenvironment{probtag}[1]{\begin{prob}}{\end{prob}}
\newenvironment{explain}{}{}
\newenvironment{explain*}{}{}
\newenvironment{editorial}{}{}
\newenvironment{digress}{}{}
\newenvironment{digress*}{}{}
\newenvironment{history}{}{}
\newenvironment{reference}{}{}
\newcommand*{\olfileid}[3]{}
\newcommand*{\olsection}[1]{\section*{#1}}
\newcommand*{\olchapter}[3]{}
\newcommand*{\olpart}[2]{}
\newcommand*{\olimport}[1]{}
\newcommand*{\olinclude}[1]{}
\newcommand*{\olpath}[1]{}
\newcommand*{\ollangid}[1]{}
\newcommand*{\problemtag}[1]{}
\newcommand*{\problemauthor}[1]{}
\newcommand*{\ollabel}[1]{\label{#1}}
% \olref takes 0-3 OPTIONAL arguments plus one mandatory label:
%     \olref{lab}  |  \olref[from]{lab}  |  \olref[from][to]{lab}  |  \olref[a][b][c]{lab}
% A single-argument definition silently swallows "[sth]" and emits a broken reference
% (measured on the host: the text came out as "??iso]x"). The chain below consumes the
% brackets and cross-references only the mandatory argument.
\makeatletter
\newcommand{\olref}{\@ifnextchar[{\olref@i}{\olref@final}}
\def\olref@i[#1]{\@ifnextchar[{\olref@ii}{\olref@final}}
\def\olref@ii[#1]{\@ifnextchar[{\olref@iii}{\olref@final}}
\def\olref@iii[#1]{\olref@final}
\def\olref@final#1{\ref{#1}}
\makeatother
\newcommand*{\olnewpage}{\clearpage}

% ===================== OLP compatibility layer (1) =====================

\makeatletter

% ---------------- proof trees (OLP's prooftree is a bussproofs derivative) ----------------
\providecommand{\AxiomC}[1]{}
\providecommand{\UnaryInfC}[1]{}
\providecommand{\BinaryInfC}[1]{}
\providecommand{\TrinaryInfC}[1]{}
\providecommand{\QuaternaryInfC}[1]{}
\providecommand{\QuinaryInfC}[1]{}
\providecommand{\RightLabel}[1]{}
\providecommand{\LeftLabel}[1]{}
\providecommand{\doubleLine}{}
\providecommand{\noLine}{}
\providecommand{\DisplayProof}{}
\providecommand{\fCenter}{}
\newenvironment{prooftree}{}{}
\newenvironment{ltprooftree}{}{}
\newenvironment{oltableau}{}{}

% ---------------- logical notation ----------------
\providecommand{\lif}{\mathbin{\rightarrow}}
\providecommand{\lforall}{\mathord{\forall}}
\providecommand{\lexists}{\mathord{\exists}}
\providecommand{\ltrue}{\top}
\providecommand{\lfalse}{\bot}
\providecommand{\Ltrue}{\top}
\providecommand{\Lfalse}{\bot}
\providecommand{\Sequent}{\mathrel{\Rightarrow}}
\providecommand{\Proves}{\mathrel{\vdash}}
\providecommand{\Entails}{\mathrel{\models}}
\providecommand{\Subst}[3]{#1[#2/#3]}
\providecommand{\sFmla}[2]{#1(#2)}
\providecommand{\pFmla}[2]{#1(#2)}
\providecommand{\formula}{\relax}
\providecommand{\True}{\top}
\providecommand{\False}{\bot}
\providecommand{\Args}[1]{#1}
\providecommand{\Ts}{\ensuremath{\top}}
\providecommand{\Fi}{\ensuremath{\bot}}
\providecommand{\lfalsum}{\bot}

% ---------------- objects, structures, metatheory ----------------
\providecommand{\Obj}{}
\providecommand{\Struct}{\ensuremath{\mathfrak{A}}}
\providecommand{\Domain}{\ensuremath{|\mathfrak{A}|}}
\providecommand{\Assign}{\ensuremath{s}}
\providecommand{\Sat}[2]{\ensuremath{#1 \models #2}}
\providecommand{\mSat}[3]{\ensuremath{#1 \models #2[#3]}}
\providecommand{\mModel}[2]{\ensuremath{#1 \models #2}}
\providecommand{\Lang}{\ensuremath{\mathcal{L}}}
\providecommand{\Th}{\ensuremath{\mathrm{Th}}}
\providecommand{\Log}{\ensuremath{\mathrm{L}}}
\providecommand{\tuple}[1]{\langle #1 \rangle}
\providecommand{\num}[1]{\underline{#1}}
\providecommand{\cfind}[1]{\varphi_{#1}}
\providecommand{\Nat}{\ensuremath{\mathbb{N}}}
\providecommand{\Mult}{\ensuremath{\mathrm{Mult}}}
\providecommand{\Add}{\ensuremath{\mathrm{Add}}}
\providecommand{\Exp}{\ensuremath{\mathrm{Exp}}}
\providecommand{\Zero}{\ensuremath{0}}
\providecommand{\Succ}{\ensuremath{\mathrm{S}}}
\providecommand{\lambd}{\ensuremath{\lambda}}
\providecommand{\bredpar}{\mathrel{\rightarrow_\beta}}
\providecommand{\bcd}[1]{\ensuremath{\ulcorner #1 \urcorner}}
\providecommand{\Set}[1]{\ensuremath{\{#1\}}}
\providecommand{\indcase}[1]{\par\noindent\emph{Case #1}.\quad}
\providecommand{\indfrm}[1]{\par\noindent\emph{Form #1}.\quad}
\providecommand{\iftag}[1]{#1}
\providecommand{\tagitem}[1]{\item[#1]}
\providecommand{\Elim}[1]{}
\providecommand{\Intro}[1]{}
\providecommand{\LeftR}[1]{}
\providecommand{\RightR}[1]{}
\providecommand{\TAss}[1]{}
\providecommand{\TRule}[1]{}

\makeatother
\providecommand{\olpath}{.}

% ===================== OLP compatibility layer (2) =====================
% --- second compatibility block: notation and shorthands used by the frozen OLP units ---
\providecommand{\CutCS}{\mathrm{Cut}}
\providecommand{\LeftSubproofLabel}{\relax}
\providecommand{\RightSubproofLabel}{\relax}
\providecommand{\maeh}{\relax}
\providecommand{\cheight}{\relax}
\providecommand{\ande}{\mathrm{and}_E}
\providecommand{\andi}{\mathrm{and}_I}
\providecommand{\ore}{\mathrm{or}_E}
\providecommand{\ori}{\mathrm{or}_I}
\providecommand{\typeof}{\mathord{:}}
\providecommand{\pValue}[1]{#1}
\providecommand{\Sent}{\mathrm{Sent}}
\providecommand{\cardplus}{\mathbin{+}}
\providecommand{\cardtimes}{\mathbin{\cdot}}
\providecommand{\cardexpo}{\mathbin{\hat{\ }}}
\providecommand{\cardnless}{\not<}
\providecommand{\cardle}{\leq}
\providecommand{\funfromto}[3]{#1\colon #2 \to #3}
\providecommand{\funimage}[1]{\operatorname{im}(#1)}
\providecommand{\fdefined}{\mathord{\downarrow}}
\providecommand{\canonord}{\operatorname{ord}}
\providecommand{\ordtimes}{\mathbin{\cdot}}
\providecommand{\ordeq}[2]{#1 = #2}
\providecommand{\ordneq}[2]{#1 \neq #2}
\providecommand{\ordtype}[1]{\operatorname{ord}(#1)}
\providecommand{\ordsucc}{\mathord{+1}}
\providecommand{\ordunion}{\cup}
\providecommand{\onesphere}{S^1}
\providecommand{\rotationsgroup}{\mathrm{SO}(2)}
\providecommand{\equivrep}[1]{[#1]}
\providecommand{\rlexless}{\prec}
\providecommand{\ZFC}{\mathrm{ZFC}}
\providecommand{\ZF}{\mathrm{ZF}}
\providecommand{\Rat}{\mathbb{Q}}
\providecommand{\Real}{\mathbb{R}}
\providecommand{\PosInt}{\mathbb{Z}^{+}}
\providecommand{\Id}{\mathrm{id}}
\providecommand{\Bin}{\mathrm{Bin}}
\providecommand{\pto}{\rightharpoonup}
\providecommand{\dom}{\operatorname{dom}}
\providecommand{\ran}{\operatorname{ran}}
\providecommand{\len}{\operatorname{len}}
\providecommand{\concat}{\mathbin{{}^\frown}}
\providecommand{\card}[1]{\lvert #1 \rvert}
\providecommand{\Pow}{\mathcal{P}}
\providecommand{\Setabs}[2]{\{\, #1 \mid #2 \,\}}
\providecommand{\comp}{\circ}
\providecommand{\Char}{\chi}
\providecommand{\Prf}{\operatorname{Prf}}
\providecommand{\Prov}{\operatorname{Prov}}
\providecommand{\Proj}{\operatorname{proj}}
\providecommand{\Gn}{\operatorname{gn}}
\providecommand{\gn}{\operatorname{gn}}
\providecommand{\Value}[1]{\operatorname{val}(#1)}
\providecommand{\Deduce}{\mathrm{D}}
\providecommand{\Dual}{\mathrm{d}}
\providecommand{\Ax}{\mathrm{Ax}}
\providecommand{\pSat}[3]{#1 \models #2\,[#3]}
\providecommand{\pAssign}{\sigma}
\providecommand{\bforall}{\forall}
\providecommand{\bexists}{\exists}
\providecommand{\liff}{\mathbin{\leftrightarrow}}
\providecommand{\FV}{\operatorname{FV}}
\providecommand{\UnaryInf}{\ensuremath{\mid}}
\providecommand{\BinaryInf}{\ensuremath{\mid}}
\providecommand{\depth}{\operatorname{depth}}
\providecommand{\oliflabeldef}[3]{#1}
\providecommand{\TMblank}{\sqcup}
\providecommand{\TMstroke}{\mid}
\providecommand{\TMstay}{\mathrm{S}}
\providecommand{\TMright}{\mathrm{R}}
\providecommand{\TMleft}{\mathrm{L}}
\providecommand{\TMendtape}{\triangleright}
\providecommand{\umin}{\mu}
\providecommand{\Sum}{\sum}
\providecommand{\citep}[1]{\cite{#1}}
\providecommand{\citet}[1]{\cite{#1}}
\providecommand{\citeauthor}[1]{\cite{#1}}
\providecommand{\citeyear}[1]{\cite{#1}}

% ===================== OLP compatibility layer (3) =====================
% --- third compatibility block: remaining notation seen in the frozen OLP units ---
\providecommand{\iso}{\cong}
\providecommand{\nszero}{\mathrm{0}}
\providecommand{\nssucc}{\mathrm{S}}
\providecommand{\nsplus}{\mathbin{+}}
\providecommand{\nssless}{<}
\providecommand{\nsless}{<}
\providecommand{\PL}{\mathrm{PL}}
\providecommand{\Nec}{\mathord{\Box}}
\providecommand{\RK}{\mathrm{RK}}
\providecommand{\ST}{\mathrm{ST}}
\providecommand{\SSubst}[3]{#1[#2/#3]}
\providecommand{\subst}[3]{#1[#2/#3]}
\providecommand{\tagfalse}{\mathrm{F}}
\providecommand{\cutr}{\mathrm{Cut}}
\providecommand{\cardneq}{\neq}
\providecommand{\issucc}[1]{\operatorname{succ}(#1)}
\providecommand{\tup}[1]{\langle #1 \rangle}
\providecommand{\Var}{\mathrm{Var}}
\providecommand{\Bew}{\operatorname{Bew}}
\providecommand{\Con}{\operatorname{Con}}
\providecommand{\Cons}{\operatorname{Cons}}
\providecommand{\pDeduce}{\mathrm{D}}
\providecommand{\psat}{\models}
\providecommand{\den}[1]{\llbracket #1 \rrbracket}
\providecommand{\Sem}[2]{\llbracket #1 \rrbracket_{#2}}
\providecommand{\DomainOf}[1]{\lvert #1 \rvert}
\providecommand{\Form}{\mathrm{Form}}
\providecommand{\Terms}{\mathrm{Term}}
\providecommand{\Atom}{\mathrm{Atom}}
\providecommand{\FVars}[1]{\operatorname{FV}(#1)}
\providecommand{\fv}[1]{\operatorname{fv}(#1)}
\providecommand{\tuplep}[1]{\langle #1 \rangle}
\providecommand{\restrict}[2]{#1 \upharpoonright #2}
\providecommand{\sto}{\rightarrow}
\providecommand{\FmlaSet}[1]{\mathrm{Fmla}_{#1}}
\providecommand{\Thm}{\mathrm{Thm}}
\providecommand{\Eval}{\operatorname{val}}
\providecommand{\blank}{\sqcup}
\providecommand{\blankTape}{\sqcup}
\providecommand{\Consis}{\operatorname{Con}}
\providecommand{\drs}{\mathrm{d}}
\providecommand{\VRule}{\mid}
\providecommand{\Bet}{\mathrm{B}}
\providecommand{\Seg}[1]{\mathrm{Seg}_{#1}}
\providecommand{\Enums}{\mathrm{En}}
\providecommand{\Z}{\mathbb{Z}}
\providecommand{\N}{\mathbb{N}}
\providecommand{\abs}[1]{\lvert #1 \rvert}
\providecommand{\floor}[1]{\lfloor #1 \rfloor}
\providecommand{\ceil}[1]{\lceil #1 \rceil}

% ===================== OLP compatibility layer (4) =====================
% --- fourth compatibility block: harvested from the per-candidate compile logs ---
\providecommand{\beredpar}{\mathrel{\rightarrow_\beta}}
\makeatletter
\def\bered{\@ifnextchar[{\bered@i}{\mathrel{\rightarrow_\beta}}}
\def\bered@i[#1]{\mathrel{\rightarrow_\beta}}
\def\bred{\@ifnextchar[{\bred@i}{\mathrel{\rightarrow_\beta}}}
\def\bred@i[#1]{\mathrel{\rightarrow_\beta}}
\makeatother
\providecommand{\beredone}{\mathrel{=_\beta}}
\providecommand{\bredone}{\mathrel{=_\beta}}
\providecommand{\xredone}{\mathrel{=_\beta}}
\providecommand{\aconvone}{\mathrel{=_\alpha}}
\providecommand{\cardless}{<}
\providecommand{\tagrefs}[1]{#1}
\providecommand{\ordplus}{\mathbin{+}}
\providecommand{\funrestrictionto}[2]{#1 \upharpoonright #2}
\providecommand{\fundefined}{\mathord{\uparrow}}
\providecommand{\Trm}{\mathrm{Trm}}
\providecommand{\OProv}{\operatorname{Prov}}
\providecommand{\OPrf}{\operatorname{Prf}}
\providecommand{\aeq}{\mathrel{=}}
\providecommand{\mClass}[1]{[#1]}
\providecommand{\equivclass}[1]{[#1]}
\providecommand{\pheight}{\operatorname{ht}}
\providecommand{\phy}{\operatorname{ht}}
\providecommand{\Cut}{\mathrm{Cut}}
\providecommand{\QuantRank}{\operatorname{qr}}
\providecommand{\elemequiv}{\mathrel{\equiv}}
\providecommand{\PIso}{\mathrm{PIso}}
\providecommand{\Theory}{\mathrm{Th}}
\providecommand{\supstrict}{\subsetneq}
\providecommand{\trcl}{\operatorname{trcl}}
\providecommand{\cref}[1]{\ref{#1}}
\providecommand{\Discharge}{\mathrm{Dis}}
\providecommand{\Weakening}{\mathrm{Wk}}
\providecommand{\ZFminus}{\mathrm{ZF}^{-}}
\providecommand{\RK}{\mathrm{RK}}
\makeatletter
\renewcommand{\Dual}{\@ifnextchar[{\Dual@i}{\ifmmode\mathrm{d}\else$\mathrm{d}$\fi}}
\def\Dual@i[#1]{\ifmmode\mathrm{d}[#1]\else$\mathrm{d}[#1]$\fi}
\makeatother
\providecommand{\RightR}{\mathrm{R}}
\providecommand{\LeftR}{\mathrm{L}}
\providecommand{\Intro}{\mathrm{I}}
\providecommand{\Axiom}{\mathrm{Ax}}
\providecommand{\formula}[1]{#1}
\providecommand{\stared}{\ast}
\providecommand{\oset}[2]{\{\, #1 \mid #2 \,\}}
\providecommand{\setof}[2]{\{\, #1 \mid #2 \,\}}
\providecommand{\pow}[1]{\mathcal{P}(#1)}
\providecommand{\ordpair}[2]{\langle #1, #2 \rangle}
\providecommand{\ordtriple}[3]{\langle #1, #2, #3 \rangle}
\providecommand{\union}{\cup}
\providecommand{\inter}{\cap}
\providecommand{\seq}[1]{\langle #1 \rangle}
\providecommand{\SuchThat}{\mid}
\providecommand{\where}{\mid}
\providecommand{\TVar}{\mathrm{TVar}}
\providecommand{\Bool}{\mathrm{Bool}}
\providecommand{\Val}{\operatorname{Val}}
\providecommand{\FV}{\operatorname{FV}}
\providecommand{\TL}[1]{\mathrm{TL}}
\providecommand{\xstared}{\ast}

% ===================== OLP compatibility layer (5) =====================
% --- fifth compatibility block: rescues the quarantined OLP candidates ---
\providecommand{\ident}{\operatorname{id}}
\providecommand{\setrank}{\operatorname{rank}}
\providecommand{\cardfont}{\mathfrak{c}}
\providecommand{\Complement}{\overline}
\providecommand{\QR}{\mathrm{QR}}
\providecommand{\becd}{\ensuremath{\ulcorner}}
\providecommand{\disjointsum}{\mathbin{\oplus}}
\providecommand{\cardeq}[2]{#1 = #2}
\providecommand{\xred}{\mathrel{\rightarrow}}
\providecommand{\cardless}{<}
\providecommand{\cardplus}{\mathbin{+}}
\providecommand{\cardtimes}{\mathbin{\cdot}}
\providecommand{\cardexpo}{\mathbin{\hat{\ }}}
\providecommand{\cardneq}{\neq}
\providecommand{\cardle}{\leq}
\providecommand{\cardnless}{\not<}
% --- sixth compatibility block: names reported by real compile logs ---
\providecommand{\ORefut}{\operatorname{Refut}}
\providecommand{\ORProv}{\operatorname{Prov}}
\providecommand{\cutrank}{\operatorname{cr}}
\providecommand{\Refut}{\operatorname{Refut}}
\providecommand{\DischargeRule}{\mathrm{Dis}}
\providecommand{\Contraction}{\mathrm{C}}
\providecommand{\bottomAlignProof}{\relax}
\providecommand{\FalseCl}{\bot}
\providecommand{\red}{\mathrel{\rightarrow}}
\providecommand{\varAssign}{\sigma}
\providecommand{\defiff}{\mathrel{\leftrightarrow}}
\providecommand{\FalseInt}{\bot}
\providecommand{\fn}{\operatorname{fn}}
\providecommand{\stared}{\ast}
\providecommand{\dotted}{\cdot}
\providecommand{\Gamma}{\Gamma}
% --- sixth compatibility block: names reported by real compile logs ---
\providecommand{\ORefut}{\operatorname{Refut}}
\providecommand{\ORProv}{\operatorname{Prov}}
\providecommand{\cutrank}{\operatorname{cr}}
\providecommand{\Refut}{\operatorname{Refut}}
\providecommand{\DischargeRule}{\mathrm{Dis}}
\providecommand{\Contraction}{\mathrm{C}}
\providecommand{\bottomAlignProof}{\relax}
\providecommand{\FalseCl}{\bot}
\providecommand{\red}{\mathrel{\rightarrow}}
\providecommand{\varAssign}{\sigma}
\providecommand{\defiff}{\mathrel{\leftrightarrow}}
\providecommand{\FalseInt}{\bot}
\providecommand{\fn}{\operatorname{fn}}
\providecommand{\stared}{\ast}
\providecommand{\dotted}{\cdot}
\providecommand{\Gamma}{\Gamma}

% --- tag-macro semantic fix (2026-10-13, final): \iftag CHOOSES a branch and the OLP proof body is
% always the SECOND group (\iftag{tag}{body}, {tag}{body}{c}, {tag}{body}{c}{d}); a definition that
% prints the tag literal drops the proof body. Bodies contain blank lines, hence \long.
\makeatletter
\providecommand{\iftag}{}
\long\def\iftag#1#2{\@ifnextchar\bgroup{\iftag@three{#2}}{\iftag@stop{#2}}}
\long\def\iftag@three#1#2{\iftag@stop{#1}\@ifnextchar\bgroup{\iftag@eat}{\relax}}
\long\def\iftag@eat#1{}
\long\def\iftag@stop#1{#1}
\providecommand{\oliflabeldef}[2]{#1}
\long\def\oliflabeldef#1#2{\@ifnextchar\bgroup{\@gobble}{\iftag@stop{#2}}}
\providecommand{\tagrefs}[1]{}
\providecommand{\settag}[1]{}
\makeatother
\begin{document}

% ==== OLP file: content/proof-theory/propositions-as-types/introduction.tex ====
% part: intuitionistic-logic
% chapter: propositions-as-types
% section: introduction



\olfileid{int}{pty}{int}

\olsection{Introduction}

The lambda calculus is a simple model of computation. It operates on
terms built from variables. If $N$ and $M$ are terms, then $(NM)$ is
also a term (``$N$ applied to~$M$''). If $N$ is a term, then
$\lambd[x][N]$ is a term. If $N$ contains the variable~$x$, then the
corresponding occurrences of~$x$ in $\lambd[x][N]$ are bound by the
$\lambd[x]$ operator, and so are no longer free. A term of the form
$(\lambd[x][N]M)$ is said to $\beta$-contract to $\Subst{N}{M}{x}$
(the result of replacing every free occurrence of $x$ in~$N$ by $M$).
A term $N$ $\beta$-converts to~$N'$ if $N'$ results from $N$ by
$\beta$-contracting a subterm, and we write $N \redone N'$. A
computation in the $\lambd$-calculus is a sequence $N \redone N_1
\redone N_2 \redone \dots$. If this sequence ends in a term $N_k$
where no subterm can be $\beta$-contracted, it is called \emph{normal}
or a normal form of~$N$. Think of computation in the $\lambd$-calculus
as such sequences. The computation halts if it results in a normal
form, and the normal form is the result of the computation. This
simple model of computations turns out to be Turing complete: every
partial recursive function can be computated by a term in the $\lambd$-calculus.

Haskell Curry and William Howard independently discovered a close
similarity between natural deduction and the $\lambd$-calculus.
Intuitively, a term $\lambd[x][N]$ is a function: $x$ is the argument;
$N$ tells us how to compute the value given~$x$. A term
$(\lambd[x][N]M)$ then is applying the function $\lambd[x][N]$ to an
argument~$M$, and $\beta$-contraction tells us how to evaluate it:
replace $x$ in $N$ by~$M$ (and keep evaluating the resulting term).
Now think of these as functions with a fixed domain and range. If the
domain of $\lambd[x][N]$ is $A$ and and its range~$B$, then we should
think of the variable~$x$ as only taking arguments from~$A$ and $N$ as
producing values from~$B$. Consequently, $(\lambd[x][N]M)$ only should
make sense if $\lambd[x][N]$ is a function $A \to B$, and $M \in A$.
If we want to add this information to the $\lambd$-calculus we obtain
the \emph{typed} $\lambd$-calculus. In the typed $\lambd$-calculus not
every expression $(\lambd[x][N]M)$ is legal, but only those where the
``types'' of $\lambd[x][N]$ and $M$ fit together, i.e., $\lambd[x][N]$
has type $A \lif B$ and $M$ has type~$A$. The type of $\lambd[x][N]$
is $A \to B$ only if the type of $x$ is $A$ and the type of $N$
is~$B$. In general, not every expression $(M_1M_2)$ is legal. Only
those terms $(M_1M_2)$ are legal where $M_1$ has type $A \to B$ and
$M_2$ has type~$A$; if so, then $(M_1M_2)$ has type~$B$. 

These typing rules now clearly correspond to !!{derivation}s in
natural deduction, where the types are represented by !!{formula}s,
and the !!{derivation}s themselves to $\lambd$-terms. For instance,
!!a{derivation} of $!A \lif !B$ corresponds to a term
$\lambd[\typeof{x}{!A}][N]$, which has the function type
$!A \lif !B$. The inference rules of natural deduction correspond to
typing rules in the typed lambda calculus, e.g.,
\begin{prooftree}
  \AxiomC{$\Discharge{!A}{x}$}
  \DeduceC{$!B$}
  \DischargeRule{\Intro\lif}{x}
  \UnaryInfC{$!A \lif !B$}
  \DisplayProof\bottomAlignProof
  \quad\text{corresponds to}\quad
  \Axiom$x: !A \fCenter N: !B$
  \UnaryInf$ \fCenter \lambd[\typeof{x}{!A}][N] : !A \lif !B$
\end{prooftree}
where the rule on the right means that if $x$ is of type~$!A$ and $N$
is of type~$!B$, then $\lambd[\typeof{x}{!A}][N]$ is of type~$!A \lif
!B$.

The $\Elim\lif$ rule corresponds to the typing rule for composition
terms, i.e.,
\[
  \AxiomC{$!A \lif !B$}
  \AxiomC{$!A$}
  \RightLabel{\Elim\lif}
  \BinaryInfC{$!B$}
  \DisplayProof
  \quad\text{corresponds to}\quad
  \Axiom$\fCenter M_1: !A \lif !B$
  \Axiom$\fCenter M_2: !A$
  \BinaryInf$ \fCenter (M_1M_2) : !B$
  \DisplayProof
\]
If a \Intro\lif{} rule is followed immediately by a \Elim\lif{} rule,
the !!{derivation} can be simplified:
\begin{gather*}
  \AxiomC{$\Discharge{!A}{x}$}
  \DeduceC{$!B$}
  \DischargeRule{\Intro\lif}{x}
  \UnaryInfC{$!A \lif !B$}
  \AxiomC{}
  \DeduceC{$!A$}
  \RightLabel{\Elim\lif}
  \BinaryInfC{$!B$}
  \DisplayProof
  \quad\redone\quad
  \AxiomC{}
  \DeduceC{$!A$}
  \DeduceC{$!B$}
  \DisplayProof
\end{gather*}
which corresponds to $\beta$-contraction of lambda terms
\[
(\lambd[\typeof{x}{!A}][N]M) \quad\redone\quad
\Subst{N}{M}{x}.
\]

Similar correspondences hold between the rules for $\land$ and ``product''
types, and between the rules for $\lor$ and ``sum'' types.

This correspondence between terms in the simply typed lambda calculus
and natural deduction !!{derivation}s is called the ``Curry--Howard'',
``proofs-as-programs,'' or ``propositions as types'' correspondence.
In addition to !!{formula}s (propositions) corresponding to types, and
proofs to terms, we can summarize the correspondences as follows:
\begin{center}
  \begin{tabular}{c c c}
    logic & program \\
    \hline
    proposition/!!{formula} & type \\
    proof & term \\
    assumption & variable \\
    discharged assumption & bound variable \\
    undischarged assumption & free variable \\
    $!A \lif !B$ & function type \\
    $!A \land !B$ & product type \\
    $!A \lor !B$ & sum type \\
    $\lfalse$ & empty type \\
    \hline
  \end{tabular}
\end{center}

The Curry--Howard correspondence is one of the cornerstones of
automated proof assistants and type checkers for programs, since
checking a proof witnessing a proposition (as we did above) amounts to
checking if a program (term) has the declared type.

% ==== OLP file: content/proof-theory/propositions-as-types/normalization.tex ====
% part: intuitionistic-logic
% chapter: propositions-as-types
% section: normalization



\olfileid{int}{pty}{nor}

\olsection{Normalization}

In this section we prove that, via some reduction order, any deduction
can be reduced to a normal deduction, which is called the
\emph{normalization property}. We will make use of the
propositions-as-types correspondence: we show that every proof term
can be reduced to a normal form; normalization for natural deduction
!!{derivation}s then follows.

Firstly we define some functions that measure the complexity of terms.
The \emph{length}~$\len{!A}$ of !!a{formula}s is defined by
\begin{align*}
  \len{p} &= 0\\
  \len{!A \land !B} &= \len{!A} + \len{!B} + 1\\
  \len{!A \lor !B} &= \len{!A} + \len{!B} + 1 \\
  \len{!A \lif !B} &= \len{!A} + \len{!B} + 1.
\end{align*}
The complexity of a redex~$M$ is measured by its \emph{cut
  rank}~$\cutrank{M}$:
\begin{align*}
  \cutrank{(\lambd[\typeof{x}{!A}][\typeof{N}{!B}])Q} &=
  \len{!A} + \len{!B} + 1 \\
  \cutrank{\ande{i}{\andi{\typeof{M}{!A}}{\typeof{N}{!B}}}} &=
  \len{!A} + \len{!B} + 1 \\
  \cutrank{\ore{\ori{i}{!A_{i}}{\typeof{M}{!A_i}}}
    {\typeof{x_1}{!A_1}}{\typeof{N_1}{!C}}
    {\typeof{x_2}{!A_2}}{\typeof{N_2}{!C}}} &=
  \len{!A} + \len{!B} + 1
\end{align*}
The complexity of a proof term is measured by the most complex
redex in it, and $0$ if it is normal:
\[
  \maxrank{M} = \max \{\cutrank{N} | N \text{ is a sub term of } M \text{
    and is redex}\}
\]

\begin{lem}\ollabel{lem:subst}
  If $\Subst{M}{\typeof{N}{!A}}{\typeof{x}{!A}}$ is a redex and $M
  \not\ident x$, then one of the following cases holds:
  \begin{enumerate}
  \item $M$ is itself a redex, or
  \item $M$ is of the form $\ande{i}{x}$, and $N$ is of the form
    $\andi{P_1}{P_2}$
  \item $M$ is of the form $\ore{i}{x_1}{P_1}{x_2}{P_2}$, and $N$ is
    of the form $\ori{i}{}{Q}$
  \item $M$ is of the form $x Q$, and $N$ is of the form
    $\lambd[x][P]$
  \end{enumerate}
  In the first case, $\cutrank{\Subst{M}{N}{x}} = \cutrank{M}$; in the
  other cases, $\cutrank{\Subst{M}{N}{x}} = \len{!A})$.
\end{lem}
\begin{proof}
  Proof by induction on~$M$.
  \begin{enumerate}
  \item If $M$ is a single variable $y$ and $y \not\ident x$, then
    $\Subst{y}{N}{x}$ is $y$, hence not a redex.
  \item If $M$ is of the form $\andi{N_1}{N_2}$, or $\lambd[x][N]$, or
    $\ori{i}{!A}{N}$, then $\Subst{M}{\typeof{N}{!A}}{\typeof{x}{!A}}$ is
    also of that form, and so is not a redex.
  \item If $M$ is of the form $\ande{i}{P}$, we consider two cases.
    \begin{enumerate}
      \item If $P$ is of the form $\andi{P_1}{P_2}$, then $M \ident
        \ande{i}{\andi{P_1}{P_2}}$ is a redex, and clearly
        \[
        \Subst{M}{N}{x} \ident
        \ande{i}{\andi{\Subst{P_1}{N}{x}}{\Subst{P_2}{N}{x}}}
        \]
        is also a redex. The cut ranks are equal.
      \item If $P$ is a single variable, it must be $x$ to make the
        substitution a redex, and $N$ must be of the form
        $\andi{P_1}{P_2}$. Now consider
        \[
        \Subst{M}{N}{x} \ident \Subst{\ande{i}{x}}{\andi{P_1}{P_2}}{x},
        \]
        which is $\ande{i}{\andi{P_1}{P_2}}$. Its cut rank is equal to
        $\cutrank{x}$, which is $\len{!A}$.
    \end{enumerate}
  \end{enumerate}
  The cases of $\ore{N}{x_1}{N_1}{x_2}{N_2}$ and $P Q$ are similar.
\end{proof}

\begin{lem}
  If $M$ contracts to $M'$, and $\cutrank{M} > \cutrank{N}$ for all
  proper redex sub-terms $N$ of $M$, then $\cutrank{M} >
  \maxrank{M'}$.
\end{lem}

\begin{proof}
  Proof by cases.
  \begin{enumerate}
  \item If $M$ is of the form $\ande{i}{\andi{M_1}{M_2}}$, then $M'$ is
    $M_i$; since any sub-term of $M_i$ is also proper
    sub-term of $M$, the claim holds.
  \item If $M$ is of the form
    $(\lambd[\typeof{x}{!A}][N])\typeof{Q}{!A}$, then $M'$ is
    $\Subst{N}{\typeof{Q}{!A}}{\typeof{x}{!A}}$.  Consider a redex in
    $M'$.  Either there is corresponding redex in $N$ with equal cut
    rank, which is less than $\cutrank{M}$ by assumption, or the cut
    rank equals $\len{!A}$, which by definition is less than
    $\cutrank{(\lambd[x^{!A}][N])Q}$.
  \item If $M$ is of the form
    \[
    \ore{\ori{i}{}{\typeof{N}{!A_i}}}
        {\typeof{x_1}{!A_1}}{\typeof{N_1}{!C}}
        {\typeof{x_2}{!A_2}}{\typeof{N_2}{!C}},
    \]
    then $M' \ident \Subst{N_i}{N}{\typeof{x_i}{!A_i}}$.  Consider a
    redex in~$M'$. Either there is corresponding redex in $N_i$ with
    equal cut rank, which is less than $\cutrank{M}$ by assumption; or
    the cut rank equals $\len{!A_i}$, which by definition is less than
    $\cutrank{\ore{\ori{i}{}{\typeof{N}{!A_i}}}
      {\typeof{x_1}{!A_1}}{\typeof{N_1}{!C}}
      {\typeof{x_2}{!A_2}}{\typeof{N_2}{!C}}}$.
  \end{enumerate}
\end{proof}

\begin{thm}
  All proof terms reduce to normal form; all !!{derivation}s reduce to
  normal !!{derivation}s.
\end{thm}

\begin{proof}
  The second follows from the first. We prove the first by complete
  induction on $m=\maxrank{M}$, where $M$ is a proof term.
  \begin{enumerate}

  \item If $m=0$, $M$ is already normal.
  \item
    Otherwise, we proceed by induction on~$n$, the number of
    redexes in $M$ with cut rank equal to~$m$.
    \begin{enumerate}
    \item If $n=1$, select any redex $N$ such that $m = \cutrank{N} >
      \cutrank{P}$ for any proper sub-term $P$ which is also a redex
      of course. Such a redex must exist, since any term only has
      finitely many subterms.

      Let $N'$ denote the reductum of $N$. Now by the lemma
      $\maxrank{N'} < \maxrank{N}$, thus we can see that $n$, the
      number of redexes with $\cutrank=m$ is decreased. So $m$ is
      decreased (by $1$ or more), and we can apply the inductive
      hypothesis for~$m$.
    \item For the induction step, assume $n>1$. the process is
      similar, except that $n$ is only decreased to a positive number
      and thus $m$ does not change. We simply apply the induction
      hypothesis for~$n$.
    \end{enumerate}
\end{enumerate}
\end{proof}

The fact that typed $\lambd$-terms normalize means that every such
term \emph{has a normal form}. If we think of $\beta$-reducing
$\lambd$-terms to their normal form as executing a computation and the
normal form as the value of the computation, this means that every
computation in the typed $\lambd$-calculus eventually halts an
produces a value. Moreover, type preservation ensures that the value
has the right type. E.g., if $N$ has type $!A \lif !B$ (i.e., it is a
function that is supposed to work for arguments of type~$!A$ and
produces values of type~$!B$) and we apply it to a term~$M$ of type~$!A$
(the input), $(NM)$ reduces to a term $N'$ (the output) and that
output is guaranteed to be of type~$!B$.

In fact, terms in the typed $\lambd$-calculus don't just reduce to
normal forms for one way of applying $\beta$-reductions to subterms,
but \emph{any} way of applying $\beta$-reductions to subtems
eventually terminates in a normal form. This property is called
\emph{strong normalization}. It means that not only is there a way to
guarantee that a computation eventually halts and produces a value,
but \emph{any} way to execute the computation eventually produces a
value.

How do we know that different ways of executing a computation produce
the \emph{same} value? So far we just know (from type preservation)
that all values it produces are of the same type. The typed
$\lambd$-calculus has another important property: it is
\emph{confluent} or \emph{Church-Rosser}. If $N \red N_1$ and $N \red
N_2$ then there is some term~$N'$ such that $N_1 \red N'$ and $N_2 \red
N'$. Recall that $N \red N'$ means: $N'$ results from zero or more
$\redone$-reduction steps. If $N_1$ is in normal form then the only
term $N'$ such that $N_1 \red N'$ is $N_1$ itself (via zero
$\redone$-reduction steps). It follows that if $N_1$ and $N_2$ are
both in normal form, then $N_1 = N' = N_2$. In other words, since
$\red$ is strongly normalizing, every way of reducing a term
eventually results in a normal form. Since $\red$ is Church-Rosser,
any two normal forms are equal. Hence, computation in the typed
$\lambd$-calculus always terminates and the value (normal form) is the
same regardless of how the computation proceeds.

Proving that a relation is confluent is typically difficult. However,
we can easily establish the following weaker condition:

\begin{prop}[Weak Church-Rosser Property]
If $N \redone N_1$ and $N \redone N_2$ then there is a term $N'$ such
that $N_1 \red N'$ and $N_2 \red N'$.
\end{prop}

\begin{proof}
  The two terms $N_1$ and $N_2$ result from $N$ by replacing the redex
  $M_1$ or $M_2$ by its reductum, respectively. The redexes $M_1$ and
  $M_2$ are subterms of $N$, and so cannot overlap. This means that
  either one is a subterm of the other, or they occur in different
  non-overlapping places in $N$. Consider the latter case: $N$ is of
  the form $N(M_1,M_2)$. Replacing $M_1$ by its reductum $M_1'$
  results in $N_1$ which is $N(M_1',M_2)$, and replacing $M_2$ by its
  reductum $M_2'$ results in $N_2$ which is $N(M_1,M_2')$. Both reduce
  to $N(M_1',M_2')$. For the other case, suppose $M_2$ is a sub-term
  of~$M_1$ (the other case is symmetric) and distinguish cases
  according to what kind of redex $M_1$ is. E.g., if $M_1 \ident
  \proj{1}{\pair{O_1}{O_2}}$ it reduces to $O_1$. If $M_2$ is a
  subterm of $O_1$, then converting $M_2$ results in $O_1'$. So
  converting $M_1$ first and then $M_2$ results in $O_1'$. But
  converting $M_2$ first results in $\proj{1}{\pair{O_1'}{O_2}}$ and
  this reduces also to~$O_2$.
\end{proof}

\begin{lem}[Newman's Lemma]
If $\red$ is strongly normalizing and has the weak Church-Rosser
property, it is confluent.
\end{lem}

\begin{proof}
  Since $\red$ is strongly normalizing, any reduction sequence results
  in a normal form. Normal forms are unique iff $\red$ is confluent.
  So suppose some term $N$ has two distinct normal forms $N'$ and
  $N''$, the result of reduction sequences
  \begin{align*}
  N & = N_1' \redone N_2' \redone \dots \redone N_k' = N' \text{ and}\\
  N & = N_1'' \redone N_2'' \redone \dots \redone N_l'' = N''
  \end{align*}
  (The reduction sequence must have length $>0$ since otherwise $N$
  would already be in normal form.)

  If $N_1' = N_1''$ then it has two distinct normal forms ($N'$ and
  $N''$). If $N_1' \neq N_1''$ then by weak Church-Rosser there is
  some $N'''$ such that $N_1' \red N'''$ and $N_1'' \red N'''$. Since
  $N' \neq N''$, either $N''' \neq N'$ or $N''' \neq N''$.
  Consequently, either $N_1'$ or $N_1''$ has two distinct normal
  forms. We've shown that if $N$ has two distinct normal forms, there
  is an $N_1$ such that $N \redone N_1$ and $N_1$ has two distinct
  normal forms. Obviously we can keep going and produce $N \redone N_1
  \redone N_2 \redone \dots$, but this is impossible because $\redone$
  is strongly normalizing.
\end{proof}



% ==== OLP file: content/proof-theory/propositions-as-types/proof-terms.tex ====
% part: intuitionistic-logic
% chapter: propositions-as-types
% section: proof-terms



\olfileid{int}{pty}{ter}

\olsection{Proof Terms}

We label assumptions in \Log{N1i} with variables (e.g., $!A^x$) as
discharge labels. In \Log{N2i}, every sequent $\Gamma \Sequent !A$
has a list of variable--!!{formula} pairs in the context~$\Gamma$ on
the left. Thus, in \Log{N2i}, every sequent in !!a{proof} carries not
only the information what the !!{proof} so far has !!{prove}d (i.e.,
$!A$), but also which assumptions are open at each point (i.e.,
$\Gamma$). What if we also included in each sequent the information
\emph{how} $!A$ was !!{prove}d. To this end, we can define a
\emph{term-labelled} system \Log{tN2i}, in which sequents have the form
$\Gamma \Sequent N : !A$. $\Gamma$~and $!A$ are as before: a set of
pairs of the form $x: !B$ and the !!{formula} the sub-!!{proof} ending
in the sequent has shown to follow from the assumptions in~$\Gamma$.
$N$~will be a certain term that encodes the !!{proof} ending in the
sequent.

The terms we assign to each sequent are built from variables $x$, $y$,
\dots, using function symbols which encode the inferences carried out
so far. We need a separate function symbol for each inference rule.
The function symbol corresponding to !!{introduction} rules are called
``constructors'' and those corresponding to !!{elimination} rules
``destructors.'' Each function symbol takes as many arguments as the
rule has premises. For instance, the function symbols for
\Intro{\land} and \Elim{\land}[1] might be $c^\land$ (2-place) and
$d^\land_1$ (1-place). In \Log{tN2i} these rules then become
\[
\Axiom$\Gamma_1 \fCenter N: !A$
\Axiom$\Gamma_2 \fCenter M: !B$
\RightLabel{\Intro{\land}}
\BinaryInf$\Gamma_1, \Gamma_2 \fCenter c^\land(N, m) : !A \land !B $
\DisplayProof
\quad
\Axiom$\Gamma \fCenter N: !A \land !B$
\RightLabel{\Elim{\land}[1]}
\UnaryInf$\Gamma \fCenter d^\land_1(N) : !A$
\DisplayProof
\]
If an inference rule discharges an assumption, i.e., the rule has a
premise that indicates a labelled !!{formula}~$x:!A$ in a premise that
is not in the context of the conclusion, we have to indicate this in
the proof term as well. The function symbols associated with such rules
\emph{bind} the corresponding variables. For instance, the
\Intro{\lif} rule goes from the premise $x:!A, \Gamma \Sequent !B$ to
$\Gamma \Sequent !A \lif !B$. The constructor $c^{\lif}$ takes one
argument, but also must indicate that the variable $x$ becomes bound.
This is the proof term's way of expressing that the assumption with
label~$x$ has been discharged. For instance, if the proof term for the
premise is~$N$, we could write the proof term for the conclusion as
$c^{\lif}([x]N)$ and the rule as
\[
\Axiom$x:!A, \Gamma \fCenter N: !B$
\RightLabel{\Intro{\lif}}
\UnaryInf$\Gamma \fCenter c^{\lif}([x]N) : !A \lif !B$
\DisplayProof
\]

In order to be able to fully reconstruct !!a{proof}, we need some
additional information. If the premise !!{formula}s of a rule don't
fully determine what the conclusion !!{formula} is, we have to add
that information to the term. This is the case in \Intro{\lif}, so we
will write $c^{\lif}([x:!A]N)$. It is also the case in \Intro{\lor}
and~\FalseInt, so we use function symbols that carry that information,
e.g.,
\[
\Axiom$\Gamma \fCenter N:!A$
\RightLabel{\Intro{\lor}[1]}
\UnaryInf$\Gamma \fCenter c^{\lor{!B}}_1:!A \lor !B$
\DisplayProof
\quad
\Axiom$\Gamma \fCenter N:\lfalse$
\RightLabel{\FalseInt}
\UnaryInf$\Gamma \fCenter d^\lfalse_{!A}(N):!A$
\DisplayProof
\]

With these ideas it is possible to construct a sequent-style natural
deduction system in which every sequent is of the form $\Gamma
\Sequent N : !A$, where $N$~is a \emph{proof term}.

It turns out that there is a close connection between such proof terms
and terms in the so-called \emph{typed lambda calculus}. To respect
this connection, we'll use symbols used in the literature on the
lambda calculus rather than the generic constructor and destructor
symbols above. E.g., instead of writing $c^{\lif}([x:!A]N)$ we write
$\lambd[\typeof{x}{!A}][N]$, instead of $d^{\lif}(N,M)$ we write $(NM)$,
$\pair{N}{M}$ instead of $c^\land(N,M)$, and $\proj{i}{N}$ instead
of~$d^\land_i(N)$.

\begin{defn}
  Proof terms are inductively generated by the following rules:
  \begin{enumerate}
  \item Every variable $x$ is a proof term (Axioms).
  \item If $x$ is a variable, $!A$ is !!a{formula}, and $N$ is a proof
    term, then $\lambd[\typeof{x}{!A}][N]$ is also a proof term~(\Intro\lif).
  \item If $N$ and $M$ are proof terms, then $(NM)$ is also a proof
    term~(\Elim\lif).
  \item If $N$ and $M$ are proof terms, then $\pair{N}{M}$ is a proof
    term~(\Intro\land).
  \item If $N$ is a proof term, then $\proj{i}{N}$ is also a proof
    term, where $i$ is $1$ or~$2$~(\Elim\land[i]).
  \item If $N$ is a proof term, and $!A$ is a formula, then
    $\inj[!A]{i}{!M}$ is a proof term, where $i$ is $1$ or~$2$~(\Intro\lor[i]).
  \item If $M$, $N_1$, $N_2$ are proof terms, and $x_1$, $x_2$ are
    variables, then $\dcase{M}{x_1}{N_1}{x_2}{N_2}$ is a proof
    term~(\Elim\lor).
  \item If $N$ is a proof term and $!A$ is !!a{formula}, then
    $\abort{!A}{N}$ is proof term~(\FalseInt).
  \end{enumerate}
\end{defn}

Not every proof term is the proof term of !!a{proof}. For instance,
a term $\pair{N}{M}$ encodes a proof ending in a~\Intro{\land}
inference, so its end-!!{formula} is a conjunction~$!A \land !B$. This
can't be the major premise of a~\Elim{\lif} rule, so $(\pair{N}{M}P)$
can't be the proof term of a correct proof. Only those proof terms
that follow the rules of \Log{tN2ip} are \emph{correct} proof terms.
These rules are given in \olref{tab:tN2ip}.

\subfile{rules-tN2}


% ==== OLP file: content/proof-theory/propositions-as-types/proofs-to-terms.tex ====
% part: intuitionistic-logic
% chapter: propositions-as-types
% section: proofs-to-terms



\olfileid{int}{pty}{pt}

\olsection{Converting \usetoken{P}{derivation} to Proof Terms}

The process of converting natural deduction !!{proof}s to proof terms,
i.e., translating !!{proof}s in~\Log{N2i} to !!{proof}s in~\Log{tN2i},
is extremely straightforward. It simply amounts to writing a proof
term to the left of each !!{formula} on the right of a sequent in the
!!{derivation}, resulting in sequents of the form $\Gamma \Sequent M :
!A$.  We'll then say that $M$ \emph{witnesses}~$!A$ in
context~$\Gamma$. Which proof term to write is entirely determined by
the rules of \Log{tN2i}.

\begin{prop}
Every !!{proof}~$\delta$ of $\Gamma \Sequent !A$ in \Log{N2i}
corresponds to !!a{proof}~$\delta'$ of $\Gamma \Sequent N : !A$ in
\Log{tN2i}. The proof term $N$ is uniquely determined by~$\delta$.
\end{prop}

\begin{proof}
We proceed by induction on the height of~$\delta$. If $\delta$ is the
axiom $x:!A \Sequent !A$, then $\delta'$ is the axiom~$x:!A \Sequent
x:!A$. If $\delta$ ends in an inference rule, the inductive hypothesis
provides proof terms for each premise, and \Log{tN2i}-!!{proof}s of
sequents including these proof terms  corresponding to each premise.
We apply the corresponding rule of \Log{tN2i}; the corresponding proof
term is determined by the \Log{tN2i} inference rule.
\end{proof}

\begin{ex}
  Consider the very simple !!{proof} in \Log{N1i} of $(!A \land !B)
  \lif (!B \land !A)$:
  \[
  \AxiomC{$\Discharge{!A\land !B}{x}$}
  \RightLabel{\Elim\land[2]}
  \UnaryInfC{$!B$}
  \AxiomC{$\Discharge{!A\land !B}{x}$}
  \RightLabel{\Elim\land[1]}
  \UnaryInfC{$!A$}
  \RightLabel{\Intro\land}
  \BinaryInfC{$!B \land !A$}
  \DischargeRule{\Intro\lif}{x}
  \UnaryInfC{$(!A \land !B) \lif (!B \land !A)$}
  \DisplayProof
  \]
  In \Log{N2i} it looks like this:
  \[
  \Axiom$x : !A\land !B \fCenter !A\land !B$
  \RightLabel{\Elim\land[2]}
  \UnaryInf$x : !A\land !B \fCenter !B$
  \Axiom$x : !A\land !B \fCenter !A\land !B$
  \RightLabel{\Elim\land[1]}
  \UnaryInf$x : !A\land !B \fCenter !A$
  \RightLabel{\Intro\land}
  \BinaryInf$x : !A\land !B \fCenter !B \land !A$
  \RightLabel{\Intro\lif}
  \UnaryInf$\fCenter (!A \land !B) \lif (!B \land !A)$
  \DisplayProof
  \]
  We assign proof terms from the top down. The proof terms for the
  axioms are just the same as the variable in the context,~$x$. The
  proof terms associated with $\Elim\land[i]$ are then determined as
  $\proj{2}{x}$ and~$\proj{1}{x}$. When we apply \Intro{\land}, we
  combine these two proof terms: $\pair{\proj{2}{x}}{\proj{1}{x}}$.
  Finally, the \Intro{\lif} rule applies a $\lambd$-abstraction,
  binding~$x$ and recording that $x$ labelled the assumption~$!A \land
  !B$: $\lambd[\typeof{x}{!A \land
  !B}][\pair{\proj{2}{x}}{\proj{1}{x}}]$.
  The !!{proof} in \Log{tN2i} then is:
  \[
  \Axiom$x : !A\land !B \fCenter x: !A\land !B$
  \RightLabel{\Elim\land[2]}
  \UnaryInf$x : !A\land !B \fCenter \proj{2}{x} : !B$
  \Axiom$x : !A\land !B \fCenter x: !A\land !B$
  \RightLabel{\Elim\land[1]}
  \UnaryInf$x : !A\land !B \fCenter \proj{1}{x} : !A$
  \RightLabel{\Intro\land}
  \BinaryInf$x : !A\land !B \fCenter \pair{\proj{2}{x}}{\proj{1}{x}} : 
    !B \land !A$
  \RightLabel{\Intro\lif}
  \UnaryInf$\fCenter \lambd[\typeof{x}{!A \land
  !B}][\pair{\proj{2}{x}}{\proj{1}{x}}] : (!A \land !B) \lif (!B \land !A)$
  \DisplayProof
  \]
\end{ex}



% ==== OLP file: content/proof-theory/propositions-as-types/propositions-as-types.tex ====
% part: intuitionistic-logic
% chapter: propositions-as-types



\olchapter{int}{pty}{Propositions as Types}

\begin{editorial}
  This is a \emph{very experimental} draft of a chapter on the
  Curry--Howard correspondence.  It needs more explanation and
  motivation, and there are probably errors and omissions. The proof
  of normalization should be reviewed and expanded. There are no
  examples for the product type. Permutation and simplification
  conversions are not covered. It will make a lot more sense once
  there is also material on the (typed) lambda calculus which is
  basically presupposed here. Use with extreme caution.
\end{editorial}

\olimport{introduction}
%\olimport{sequent-natural-deduction}
\olimport{proof-terms}
\olimport{proofs-to-terms}
\olimport{terms-to-proofs}
\olimport{types}
\olimport{reduction}
\olimport{type-preservation}
\olimport{normalization}

\OLEndChapterHook


% ==== OLP file: content/proof-theory/propositions-as-types/reduction.tex ====
% part: intuitionistic-logic
% chapter: propositions-as-types
% section: reduction



\olfileid{int}{pty}{red}

\olsection{Reduction}

In natural deduction !!{derivation}s, !!a{introduction} rule that is
followed by !!a{elimination} rule is a detour that can be ``reduced''.
For instance, an $\Intro{\land}$ followed by $\Elim{\land}$ ``cancel
out,'' since the conclusion of $\Elim{\land}$ is always one of the
conjuncts of the conjunction that \Intro{\land} produces, and the
premises of $\Intro{\land}$ are both conjuncts. Thus if we need
!!a{derivation} of either conjunct, we can simply use that
!!{derivation} without introducing the conjunction followed by
eliminating it. We have
\begin{gather*}
  \bottomAlignProof
  \AxiomC{}
  \RightLabel{$\delta_1$}
  \DeduceC{$!A_1$}
  \AxiomC{}
  \RightLabel{$\delta_2$}
  \DeduceC{$!A_2$}
  \RightLabel{$\Intro{\land}$}
  \BinaryInfC{$!A_1 \land !A_2$}
  \RightLabel{$\Elim{\land}_i$}
  \UnaryInfC{$!A_i$}
  \DisplayProof\bottomAlignProof
  \quad
  \redone
  \quad
  \AxiomC{}
  \RightLabel{$\delta_i$}
  \DeduceC{$!A_i$}
  \DisplayProof
\end{gather*}

We can consider a similar relation of reduction on the associated
proof terms suggested by this simplification. If $N_1$ is the proof
term for~$\delta_1$ and $N_2$ the one for~$\delta_2$, then we can say
the proof term for the proof on the left reduces (in one step) to the
proof term for the proof on the right:
\[
  \proj{i}{\pair{N_1, N_2}} \redone N_i
\]
In the typed lambda calculus, this is the \emph{beta reduction} rule for
the product type.

In natural deduction, a pair of inferences such as a \Intro{\land}
folloed by \Elim{\land}, i.e., a pair that is subject to reduction, is
called a \emph{cut}. In the typed lambda calculus the term on the left
of $\redone$ is called a \emph{redex}, and the term to the right is
called the \emph{reductum}. Unlike untyped lambda calculus, where only
$(\lambd[x][N])Q$ is considered to be redex, in the typed lambda
calculus the syntax is extended to terms involving $\pair{N}{M}$,
$\proj{i}{N}$, $\inj{i}{!A}{N}$, $\dcase{N}{x_1}{M_1}{x_2}{M_2}$, and
$\abort{!A}{N}$, with corresponding redexes.

Similarly we have reduction for disjunction:
\begin{gather*}
  \bottomAlignProof
  \AxiomC{}
  \RightLabel{$\delta$}
  \DeduceC{$!A_i$}
  \RightLabel{$\Intro{\lor}$}
  \UnaryInfC{$!A_1 \lor !A_2$}
  \AxiomC{$\Discharge{!A_1}{x_1}$}
  \RightLabel{$\delta_1$}
  \DeduceC{$!C$}
  \AxiomC{$\Discharge{!A_2}{x_2}$}
  \RightLabel{$\delta_2$}
  \DeduceC{$!C$}
  \DischargeRule{$\Elim{\lor}$}{x_1,x_2}
  \TrinaryInfC{$!C$}
  \DisplayProof\bottomAlignProof
  \quad
  \redone
  \quad
  \AxiomC{}
  \RightLabel{$\delta$}
  \DeduceC{$!A_i$}
  \RightLabel{$\delta_i$}
  \DeduceC{$!C$}
  \DisplayProof
\end{gather*}
This corresponds to a reduction on proof terms:
\begin{gather*}
  \dcase{\inj[!A_i]{i}{M}}{x_1}{N_1}{x_2}{N_2} \redone \Subst{N_i}{M}{x_i}
\end{gather*}
where $M$ is the proof term of~$\delta$, and $N_i$ is the proof term
of~$\delta_i$. This is the beta reduction rule of for sum types.
Here, $\Subst{M}{N_i}{x}$ means replacing all free occurrences of the
variable~$x$ in $M$ by~$N_i$.

The reduction of the function type corresponds to removal of
a detour of $\Intro\lif$ followed by a $\Elim\lif$.
\begin{gather*}
  \bottomAlignProof
  \AxiomC{$\Discharge{!A}{x}$}
  \RightLabel{$\delta$}
  \DeduceC{$!B$}
  \DischargeRule{$\Intro{\lif}$}{x}
  \UnaryInfC{$!A \lif !B$}
  \AxiomC{}
  \RightLabel{$\delta'$}
  \DeduceC{$!A$}
  \RightLabel{$\Elim{\lif}$}
  \BinaryInfC{$!B$}
  \DisplayProof\bottomAlignProof
  \quad
  \redone
  \quad
  \AxiomC{}
  \RightLabel{$\delta'$}
  \DeduceC{$!A$}
  \RightLabel{$\delta$}
  \DeduceC{$!B$}
  \DisplayProof
\end{gather*}
For proof terms, this amounts to ordinary beta reduction:
\begin{gather*}
  (\lambd[\typeof{x}{!A}][N]M)
  \redone \Subst{N}{M}{x}
\end{gather*}

In addition to reduction conversions on !!{proof}s, we can consider
permutation conversions. These apply to (sub)!!{proof}s that end in a
\Elim{\lor} followed by another !!{elimination} rule, e.g.:
\begin{multline*}
  \bottomAlignProof
  \AxiomC{}
  \RightLabel{$\delta$}
  \DeduceC{$!A_i$}
  \RightLabel{$\Intro{\lor}$}
  \UnaryInfC{$!A_1 \lor !A_2$}
  \AxiomC{$\Discharge{!A_1}{x_1}$}
  \RightLabel{$\delta_1$}
  \DeduceC{$!C_1 \land !C_2$}
  \AxiomC{$\Discharge{!A_2}{x_2}$}
  \RightLabel{$\delta_2$}
  \DeduceC{$!C_1 \land !C_2$}
  \DischargeRule{$\Elim{\lor}$}{x_1,x_2}
  \TrinaryInfC{$!C_1 \land !C_2$}
  \RightLabel{\Elim{\land}[i]}
  \UnaryInfC{$!C_i$}
  \DisplayProof\bottomAlignProof
  \quad
  \redone
  \\
  \AxiomC{}
  \RightLabel{$\delta$}
  \DeduceC{$!A_i$}
  \RightLabel{$\Intro{\lor}$}
  \UnaryInfC{$!A_1 \lor !A_2$}
  \AxiomC{$\Discharge{!A_1}{x_1}$}
  \RightLabel{$\delta_1$}
  \DeduceC{$!C_1 \land !C_2$}
  \RightLabel{\Elim{\land}[i]}
  \UnaryInfC{$!C_i$}
  \AxiomC{$\Discharge{!A_2}{x_2}$}
  \RightLabel{$\delta_2$}
  \DeduceC{$!C_1 \land !C_2$}
  \RightLabel{\Elim{\land}[i]}
  \UnaryInfC{$!C_i$}
  \DischargeRule{$\Elim{\lor}$}{x_1,x_2}
  \TrinaryInfC{$!C_i$}
  \DisplayProof
\end{multline*}
If the proof terms of $\delta$, $\delta_1$, and $\delta_2$ are $M$,
$N_1$, and $N_2$, respectively, the corresponding reduction rule for
proof terms, i.e., $\lambd$ terms, is:
\[
  \proj{i}{\dcase{M}{x_1}{N_1}{x_2}{N_2}} \redone 
\dcase{M}{x_1}{\proj{i}{N_1}}{x_2}{\proj{i}{N_2}}
\]

Of course we can eliminate detours not just at the ends of !!{proof}s,
but also within !!{proof}s, by applying the reduction conversion to a
sub-!!{proof} ending in a detour. Similarly, $\beta$-reduction applies
to subterms of $\lambd$-terms. If $N$ is a subterm of $M$, $N \redone
N'$, we can replace the subterm $N$ in~$M$ by $N'$ to obtain $M'$. In
that case we also write $M \redone M'$. If there is a sequence $M = M_1
\redone M_2 \redone M_2 \redone \dots \redone M_k = M'$ (including
where $k=1$, i.e., $M = M'$) we write $M \red M'$.

A !!{proof} where no conversion can be applied to any sub-!!{proof}
is called \emph{normal}. Similarly, a $\lambd$-term where no
conversion can be applied (to any sub-term) is called \emph{normal}. A
normal $\lambd$-term is also called a \emph{normal form}.


% ==== OLP file: content/proof-theory/propositions-as-types/rules-tN2.tex ====
% Part: proof-theory
% Chapter: natural-deduction
% Section: rules-N2



\begin{table}
\begin{tabular}{@{}c@{}c@{}}
\multicolumn{2}{l}{\textbf{Axioms:}}\\
\multicolumn{2}{c}{$x:!A \Sequent x:!A$}\\[3ex]
\multicolumn{2}{l}{\textbf{Inference Rules:}}\\[3ex]
\multicolumn{2}{c}{\Axiom$x: !A, \Gamma \fCenter N:!B$
\RightLabel{\Intro{\lif}}
\UnaryInf$\Gamma \fCenter \lambd[\typeof{x}{!A}][N] : !A \lif !B $
\DisplayProof}
\\[4ex]
\multicolumn{2}{c}{\Axiom$\Gamma_1 \fCenter N:!A \lif !B$
\Axiom$\Gamma_2 \fCenter M:!A$
\RightLabel{\Elim{\lif}}
\BinaryInf$\Gamma_1, \Gamma_2 \fCenter (NM) : !B$
\DisplayProof}
\\[4ex]
\multicolumn{2}{c}{\Axiom$\Gamma_1 \fCenter N_1 : !A$
\Axiom$\Gamma_2 \fCenter N_2 : !B$
\RightLabel{\Intro{\land}}
\BinaryInf$\Gamma_1, \Gamma_2 \fCenter \pair{N_1}{N_2} : !A \land !B $
\DisplayProof}
\\[4ex]
\Axiom$\Gamma \fCenter N: !A \land !B$
\RightLabel{\Elim{\land}[1]}
\UnaryInf$\Gamma \fCenter \proj{1}{N} : !A$
\DisplayProof
&
\Axiom$\Gamma \fCenter N:!A \land !B$
\RightLabel{\Elim{\land}[2]}
\UnaryInf$\Gamma \fCenter \proj{2}{N} : !B$
\DisplayProof\\[3ex]
\\[4ex]
\Axiom$\Gamma \fCenter N: !A$
\RightLabel{\Intro{\lor}[1]}
\UnaryInf$\Gamma \fCenter \inj[!B]{1}{N} : !A \lor !B$
\DisplayProof
&
\Axiom$ \Gamma \fCenter N : !B$
\RightLabel{\Intro{\lor}[2]}
\UnaryInf$\Gamma \fCenter \inj[!A]{2}{N} : !A \lor !B$
\DisplayProof\\[3ex]
\multicolumn{2}{c}{
\begin{tabular}[b]{@{}l@{}}
\Axiom$\Gamma_1 \fCenter M : !A \lor !B$
\Axiom$x:!A, \Gamma_2 \fCenter N_1 : !C$
\Axiom$y:!B, \Gamma_3 \fCenter N_2: !C$
\RightLabel{\Elim{\lor}}
\TrinaryInf$\Gamma_1, \Gamma_2, \Gamma_3 \fCenter
\dcase{M}{x}{N_1}{y}{N_2} : !C$
\DisplayProof
\end{tabular}
}
\\[4ex]
\multicolumn{2}{c}{\Axiom$\Gamma \fCenter N: \lfalse$
\RightLabel{\FalseInt}
\UnaryInf$\Gamma \fCenter \abort{!A}{N} : !A$
\DisplayProof}
\end{tabular}
\caption{Rules of the propositional term-labelled intuitionistic
natural deduction system \Log{tN2ip}.}
\ollabel{tab:tN2ip}
\end{table}

% ==== OLP file: content/proof-theory/propositions-as-types/rules-tN3.tex ====
% Part: proof-theory
% Chapter: natural-deduction
% Section: rules-N2



\begin{table}
\begin{tabular}{@{}c@{}c@{}}
\multicolumn{2}{l}{\textbf{Axioms:}}\\
\multicolumn{2}{c}{$\Gamma \Sequent x:!A$ if $x:!A \in \Gamma$}\\[3ex]
\multicolumn{2}{l}{\textbf{Inference Rules:}}\\[3ex]
\multicolumn{2}{c}{\Axiom$x: !A, \Gamma \fCenter N:!B$
\RightLabel{\Intro{\lif}}
\UnaryInf$\Gamma \fCenter \lambd[\typeof{x}{!A}][N] : !A \lif !B $
\DisplayProof}
\\[4ex]
\multicolumn{2}{c}{\Axiom$\Gamma \fCenter N:!A \lif !B$
\Axiom$\Gamma \fCenter M:!A$
\RightLabel{\Elim{\lif}}
\BinaryInf$\Gamma \fCenter (NM) : !B$
\DisplayProof}
\\[4ex]
\multicolumn{2}{c}{\Axiom$\Gamma \fCenter N_1 : !A$
\Axiom$\Gamma \fCenter N_2 : !B$
\RightLabel{\Intro{\land}}
\BinaryInf$\Gamma \fCenter \pair{N_1}{N_2} : !A \land !B $
\DisplayProof}
\\[4ex]
\Axiom$\Gamma \fCenter N: !A \land !B$
\RightLabel{\Elim{\land}[1]}
\UnaryInf$\Gamma \fCenter \proj{1}{N} : !A$
\DisplayProof
&
\Axiom$\Gamma \fCenter N:!A \land !B$
\RightLabel{\Elim{\land}[2]}
\UnaryInf$\Gamma \fCenter \proj{2}{N} : !B$
\DisplayProof\\[3ex]
\\[4ex]
\Axiom$\Gamma \fCenter N: !A$
\RightLabel{\Intro{\lor}[1]}
\UnaryInf$\Gamma \fCenter \inj[!B]{1}{N} : !A \lor !B$
\DisplayProof
&
\Axiom$ \Gamma \fCenter N : !A$
\RightLabel{\Intro{\lor}[2]}
\UnaryInf$\Gamma \fCenter \inj[!B]{2}{N} : !B \lor !A$
\DisplayProof\\[3ex]
\multicolumn{2}{c}{
\begin{tabular}[b]{@{}l@{}}
\Axiom$\Gamma \fCenter M : !A \lor !B$
\Axiom$x:!A, \Gamma \fCenter N_1 : !C$
\Axiom$y:!B, \Gamma \fCenter N_2: !C$
\RightLabel{\Elim{\lor}}
\TrinaryInf$\Gamma \fCenter
\dcase{M}{x}{N_1}{y}{N_2} : !C$
\DisplayProof
\end{tabular}
}
\\[4ex]
\multicolumn{2}{c}{\Axiom$\Gamma \fCenter N: \lfalse$
\RightLabel{\FalseInt}
\UnaryInf$\Gamma \fCenter \abort{!A}{N} : !A$
\DisplayProof}
\end{tabular}
\caption{Rules of the propositional term-labelled intuitionistic
natural deduction system \Log{tN3ip} serving as a type assignment
system.}
\ollabel{tab:tN3ip}
\end{table}

% ==== OLP file: content/proof-theory/propositions-as-types/sequent-natural-deduction.tex ====
% part: intuitionistic-logic
% chapter: propositions-as-types
% section: sequent-natural-deduction



\olfileid{int}{pty}{snd}

\olsection{Sequent Natural Deduction}

Let us write $\Gamma \Sequent !A$ if there is a natural deduction
!!{derivation} with $\Gamma$ as !!{undischarged} assumptions and $!A$
as conclusion; or $\Sequent !A$ if $\Gamma$ is empty.

We write $\Gamma, !A_1, \dots, !A_n$ for $\Gamma \cup \{!A_1, \dots,
!A_n\}$, and $\Gamma, \Delta$ for $\Gamma \cup \Delta$.

Observe that when we have $\Gamma \Sequent !A \land !A$, meaning we
have !!a{derivation} with $\Gamma$ as !!{undischarged} assumptions and
$!A \land !A$ as end-!!{formula}, then by applying $\Elim{\land}$ at
the bottom, we can get !!a{derivation} with the same !!{undischarged}
assumptions and $!A$ as conclusion. In other words, if $\Gamma
\Sequent !A \land !B$, then $\Gamma \Sequent !A$.
\begin{prooftree}
  \Axiom$\Gamma \fCenter !A \land !B$
  \RightLabel{$\Elim{\land}$}
  \UnaryInf$\Gamma \fCenter !A$
  \DisplayProof\qquad\bottomAlignProof
  \Axiom$\Gamma \fCenter !A \land !B$
  \RightLabel{$\Elim{\land}$}
  \UnaryInf$\Gamma \fCenter !B$
\end{prooftree}
The label $\Elim{\land}$ hints at the relation with the rule of
the same name in natural deduction.

Likewise, suppose we have $\Gamma, !A \Sequent !B$, meaning we have
!!a{derivation} with !!{undischarged} assumptions $\Gamma, !A$ and
end-!!{formula}~$!B$. If we apply the $\Intro{\lif}$ rule, we have
!!a{derivation} with $\Gamma$ as !!{undischarged} assumptions and $!A
\lif !B$ as the end-!!{formula}, i.e., $\Gamma \Sequent !A \lif
!B$. Note how this has made the !!{discharge} of assumptions more
explicit.
\begin{prooftree}
  \Axiom $\Gamma, !A \fCenter !B$
  \RightLabel{$\Intro{\lif}$}
  \UnaryInf $\Gamma \fCenter !A \lif !B$
\end{prooftree}

We can draw conclusions from other rules in the same fashion, which is
spelled out as follows:
\begin{gather*}
  \Axiom $\Gamma \fCenter !A$
  \Axiom $\Delta \fCenter !B$
  \RightLabel{$\Intro{\land}$}
  \BinaryInf $\Gamma,\Delta \fCenter !A \land !B$
  \DisplayProof\\
  \Axiom $\Gamma \fCenter !A \land !B$
  \RightLabel{$\Elim{\land}_1$}
  \UnaryInf $\Gamma \fCenter !A$
  \DisplayProof
  \qquad
  \Axiom $\Gamma \fCenter !A \land !B$
  \RightLabel{$\Elim{\land}_2$}
  \UnaryInf $\Gamma \fCenter !B$
  \DisplayProof
  \\
  \Axiom $\Gamma \fCenter !A$
  \RightLabel{$\Intro{\lor}_1$}
  \UnaryInf $\Gamma \fCenter !A \lor !B$
  \DisplayProof\qquad
  \Axiom $\Gamma \fCenter !B$
  \RightLabel{$\Intro{\lor}_2$}
  \UnaryInf $\Gamma \fCenter !A \lor !B$
  \DisplayProof\\
  \Axiom $\Gamma \fCenter !A \lor !B$
  \Axiom $\Delta, !A \fCenter !C$
  \Axiom $\Delta', !B \fCenter !C$
  \RightLabel{$\Elim{\lor}$}
  \TrinaryInf $\Gamma, \Delta, \Delta' \fCenter !C$
  \DisplayProof
  \\
  \Axiom $\Gamma, !A \fCenter !B$
  \RightLabel{$\Intro{\lif}$}
  \UnaryInf $\Gamma \fCenter !A \lif !B$
  \DisplayProof
  \qquad
  \Axiom $\Delta \fCenter !A \lif !B$
  \Axiom $\Gamma \fCenter !A$
  \RightLabel{$\Elim{\lif}$}
  \BinaryInf $\Gamma, \Delta \fCenter !B$
  \DisplayProof
  \\
  \Axiom $\Gamma \fCenter \lfalse$
  \RightLabel{$\FalseInt$}
  \UnaryInf $\Gamma \fCenter !A$
  \DisplayProof
\end{gather*}

Any assumption by itself is !!a{derivation} of $!A$ from~$!A$, i.e.,
we always have $!A \Sequent !A$.

\begin{prooftree}
  \AxiomC{$ $}
  \UnaryInfC{$!A \fCenter !A$}
\end{prooftree}

Together, these rules can be taken as a calculus about what natural
deduction !!{derivation}s exist. They can also be taken as a
notational variant of natural deduction, in which each step records
not only the !!{formula} !!{derive}d but also the !!{undischarged}
assumptions from which it was !!{derive}d.

\begin{prooftree}
  \Axiom$!A \fCenter !A$
  \UnaryInf $!A \fCenter !A \lor (!A \lif \lfalse)$
  \Axiom $!B \fCenter !B$
  \BinaryInf$!A, !B\lif \fCenter \lfalse$
  \UnaryInf$(!B \fCenter !A \lif \lfalse$
  \UnaryInf$( !B \fCenter !A \lor (!A \lif \lfalse)$
  \Axiom$(!B \fCenter !B$
  \BinaryInf$(!B \fCenter \lfalse$
  \UnaryInf$\fCenter !B \lif \lfalse$
\end{prooftree}
where $!B$ is short for $(!A \lor (!A \lif \lfalse))\lif \lfalse$.


% ==== OLP file: content/proof-theory/propositions-as-types/terms-to-proofs.tex ====
% part: intuitionistic-logic
% chapter: propositions-as-types
% section: terms-to-proofs



\olfileid{int}{pty}{tp}

\olsection{Recovering \usetoken{P}{derivation} from Proof Terms}

Now let us consider the other direction: translating a proof term~$N$
back to !!a{proof} in \Log{N2i}. In general, it is only possible to do
this relative to a context that contains pairs $x : !A$ for every
variable~$x$ that occurs free in the given proof term. But let's start
with an example where there are no free variables, namely the term
\[
\lambd[\typeof{x}{!A \land
  !B}][\pair{\proj{2}{x}}{\proj{1}{x}}]
\]
It is of the form $\lambd[\typeof{x}{!A \land !B}][N_1]$. Since the
only rule of \Log{tN2i} that produces a term of this form
is~\Intro{\lif}, we know that the !!{proof} encoded by~$N$ has to end
with~\Intro{\lif}, i.e., it ends in
  \[
  \Axiom$x : !A\land !B \fCenter \pair{\proj{2}{x}}{\proj{1}{x}} : 
    !C$
  \RightLabel{\Intro\lif}
  \UnaryInf$\fCenter \lambd[\typeof{x}{!A \land
  !B}][\pair{\proj{2}{x}}{\proj{1}{x}}] : (!A \land !B) \lif !C$
  \DisplayProof
  \]
We don't yet know what $!C$ is, but we do know that the antecedent is
$!A \land !B$, and that the context of the premise has to include
$x:!A \land !B$.

Now the proof term associated with the premise is already determined:
$\pair{\proj{2}{x}}{\proj{1}{x}}$. The !!{proof} \emph{it} encodes
must end with \Intro{\land}. We still don't know what $!C$ is, but
because it is the conclusion of~\Intro{\land}, we know it must be a
conjunction, i.e., $!C \ident (!C_1 \land !C_2)$. We continue the
reconstruction:
  \[
  \Axiom$x : !A\land !B \fCenter \proj{2}{x} : !C_1$
  \Axiom$x : !A\land !B \fCenter \proj{1}{x} : !C_2$
  \RightLabel{\Intro\land}
  \BinaryInf$x : !A\land !B \fCenter \pair{\proj{2}{x}}{\proj{1}{x}} : 
    !C_1 \land !C_2$
  \RightLabel{\Intro\lif}
  \UnaryInf$\fCenter \lambd[\typeof{x}{!A \land
  !B}][\pair{\proj{2}{x}}{\proj{1}{x}}] : (!A \land !B) \lif (!C_1 \land !C_2)$
  \DisplayProof
  \]
The proof term $\pi_2(x)$ means that top sequent on the left was
obtained by \Elim{\land}[2], i.e., the premise is of the form $!D
\land !C_1$. On the right, we can determine that the inference leading
to $!C_2$ must be \Elim{\land}[1], and that the premise is of the form
$!C_2 \land !E$.
  \[
  \Axiom$x : !A\land !B \fCenter x: !D\land !C_1$
  \RightLabel{\Elim\land[2]}
  \UnaryInf$x : !A\land !B \fCenter \proj{2}{x} : !C_1$
  \Axiom$x : !A\land !B \fCenter x: !C_2\land !E$
  \RightLabel{\Elim\land[1]}
  \UnaryInf$x : !A\land !B \fCenter \proj{1}{x} : !C_2$
  \RightLabel{\Intro\land}
  \BinaryInf$x : !A\land !B \fCenter \pair{\proj{2}{x}}{\proj{1}{x}} : 
    !C_1 \land !C_2$
  \RightLabel{\Intro\lif}
  \UnaryInf$\fCenter \lambd[\typeof{x}{!A \land
  !B}][\pair{\proj{2}{x}}{\proj{1}{x}}] : (!A \land !B) \lif (!C_1 \land !C_2)$
  \DisplayProof
  \]
Since the proof terms of the topmost sequents are now variables, we
know that they must be axioms. This determines $!C_1$, $!C_2$, $!D$, and
$!E$: $!D \ident !A$ and $!C_1 \ident !B$ or else the left sequent is
not an axiom, and $!C_2 \ident !A$ and $!E \ident !B$, or else the
right sequent isn't. We have recovered the proof fully:
  \[
  \Axiom$x : !A\land !B \fCenter x: !A\land !B$
  \RightLabel{\Elim\land[2]}
  \UnaryInf$x : !A\land !B \fCenter \proj{2}{x} : !B$
  \Axiom$x : !A\land !B \fCenter x: !A\land !B$
  \RightLabel{\Elim\land[1]}
  \UnaryInf$x : !A\land !B \fCenter \proj{1}{x} : !A$
  \RightLabel{\Intro\land}
  \BinaryInf$x : !A\land !B \fCenter \pair{\proj{2}{x}}{\proj{1}{x}} : 
    !B \land !A$
  \RightLabel{\Intro\lif}
  \UnaryInf$\fCenter \lambd[\typeof{x}{!A \land
  !B}][\pair{\proj{2}{x}}{\proj{1}{x}}] : (!A \land !B) \lif (!B \land !A)$
  \DisplayProof
  \]



% ==== OLP file: content/proof-theory/propositions-as-types/type-preservation.tex ====
% part: intuitionistic-logic
% chapter: propositions-as-types
% section: type-preservation



\olfileid{int}{pty}{tpr}

\olsection{Type Preservation}

The reduction and permutation conversions considered can also be
defined directly for the systems \Log{tN2i} and \Log{tN3i}. The system
\Log{tN3i} defines the type of a $\lambd$-term~$N$ relative to the
context~$\Gamma$ as the !!{formula}~$!A$ such that $\Gamma \Sequent N
: !A$ has !!a{proof} in \Log{tN3i}. The fact that the
$\beta$-reduction rules for $\lambd$-terms exactly mirror the
reduction and permutation conversions of !!{proof}s has an important
consequence: the type of a $\lambd$-term relative to a
context~$\Gamma$ remains the same after $\beta$-conversion. This
property is called \emph{type preservation}.

\begin{prop}
If $N$ has type~$!A$ relative
to $\Gamma$ and $N \redone N'$, then $N'$ has type~$!A$ relative to~$\Gamma$.
\end{prop}

\begin{proof}
  By induction on~$N$ and the height of !!{proof}s~$\delta$ of $\Gamma
  \Sequent N : !A$ we can easily establish the following: If $M$ is a
  sub-term of~$N$, then $\delta$ contains a sub-!!{proof}~$\delta_1$
  ending in $\Gamma' \Sequent M : !B$. If $M$ is a redex, then a
  conversion applies to~$\delta_1$. By applying the reduction to the
  $\delta_1$ we obtain !!a{proof}~$\delta_1'$ of $\Gamma' \Sequent M'
  : !B$. By inspection of the conversions of \Log{tN3i}, $M'$ is the
  reductum of~$M$. Replacing $\delta_1$ in~$\delta$ by $\delta_1'$
  results in !!a{proof}~$\delta'$ of $\Gamma \Sequent N' : !A$, where
  $N'$ results from~$N$ by replacing the subterm~$M$ by~$M'$.

  For instance, consider the reduction conversion for~\Intro{\land}
  followed by \Elim{\land}:
  \[
  \bottomAlignProof
  \AxiomC{}
  \Deduce$\Gamma \fCenter N_1 : !B_1$
  \AxiomC{}
  \Deduce$\Gamma \fCenter N_2 : !B_2$
  \RightLabel{\Intro{\land}}
  \BinaryInf$\Gamma \fCenter \pair{N_1}{N_2} : !B_1 \land !B_2$
  \RightLabel{\Elim{\land}[i]}
  \UnaryInf$\Gamma \fCenter \proj{i}{\pair{N_1}{N_2}} : !B_i$
  \DisplayProof
  \quad\redone\quad
  \bottomAlignProof
  \AxiomC{}
  \Deduce$\Gamma \fCenter N_i : !B_i$
  \DisplayProof
  \]
  We see that the proof term of the !!{proof} on the right, i.e.,
  $\delta_1'$ is $N_i$. It is exactly the result of applying
  $\beta$-conversion to the proof term $\proj{i}{\pair{N_1}{N_2}}$ of
  the !!{proof} on the left~($\delta_1$).

  Or consider a \Intro{\lif} followed by an \Elim{\lif} inference, and
  the reduction conversion applied to it:
  \begin{multline*}
  \bottomAlignProof
  \AxiomC{$\delta_2$}
  \Deduce$x: !B_1, \Gamma \fCenter N : !B_2$
  \RightLabel{\Intro{\lif}}
  \UnaryInf$\Gamma \fCenter \lambd[\typeof{x}{!B_1}][N] : !B_1 \lif !B_2$
  \AxiomC{}
  \RightLabel{$\delta_3$}
  \Deduce$\Gamma \fCenter M : !B_1$
  \RightLabel{\Elim{\lif}}
  \BinaryInf$\Gamma \fCenter (\lambd[\typeof{x}{!B_1}][N] M) : !B_2$
  \DisplayProof
  \quad\redone\\
  \bottomAlignProof
  \AxiomC{}
  \RightLabel{$\delta_3$}
  \Deduce$\Gamma \fCenter M : !B_1$
  \RightLabel{$\Subst{\delta_2}{\delta_3}{x:!B_1}$}
  \Deduce$\Gamma \fCenter \Subst{N}{M}{x} : !B_2$
  \DisplayProof
  \end{multline*}
  Again the proof term of the !!{proof} on the right is exactly the
  result of $\beta$-reducing the proof term of the !!{proof} on the
  left. Of course one would have to carefuly verify (by induction on
  the height of~$\delta_2$) that if we replace all axioms of the form
  $\Gamma' \Sequent x:!B_1$ in~$\delta_2$ by the !!{proof}~$\delta_3$
  of $\Gamma \Sequent M:!B_1$ actually results in !!a{proof}
  $\Subst{\delta_2}{\delta_3}{x:!B_1}$ of $\Gamma \Sequent
  \Subst{N}{M}{x} : !B_2$.
\end{proof}

The tight correspondence between !!{proof}s and $\lambd$-terms has
another consequence. If $N$ is a $\lambd$-term of type $!A$ relative
to a context~$\Gamma$ and contains a redex~$M$ (i.e., it is not a
normal form), then the corresponding \Log{tN3i} proof~$\delta$ of
$\Gamma \Sequent N : !A$ contains a sub-!!{proof} (ending in $\Gamma'
\Sequent M: !B$) to which a conversion can be applied. If $N \redone
N'$ by replacing $M$ by its reductum, the corresponding result of
applying the reduction to~$\delta$ is !!a{proof} of $\Gamma \Sequent
N' : !A$. Hence every $\lambd$-term not in normal form $\beta$-reduces
to another. This property is called \emph{progress}.

Progress together with type preservation ensure that $\beta$-reduction
of $\lambd$-terms never ``gets stuck'': if a term is well-typed and
isn't in normal form, it reduces to another well-typed term (and that
term is of the same type).


% ==== OLP file: content/proof-theory/propositions-as-types/types.tex ====
% part: intuitionistic-logic
% chapter: propositions-as-types
% section: types



\olfileid{int}{pty}{typ}

\olsection{Types}

A proof term such as $(MN)$ by itself does not uniquely determine the
proof it is a proof term of. That's because the free variables, i.e.,
the proof terms of axioms, do not specify the !!{formula}s in the
corresponding axioms. However, if we specify these in a
context~$\Gamma$, the proof \emph{is} uniquely determined.

\begin{defn}
A \emph{context} for a proof term~$N$ is a set $\Gamma$ such that for
every free variable of~$N$, $\Gamma$ contains exactly one pair $x : !A$.
\end{defn}

Note that the definition does not require that $\Gamma$ contains
$x:!A$ \emph{only} for free variables of~$N$. This, e.g., $\Gamma =
\{x:!A, y:!B, z:!C\}$ is a context for $\pair{x}{y}$.

\begin{defn}\ollabel{defn:type}
  Suppose $\Gamma$ is a context for~$N$. The \emph{type} of~$N$
  (relative to~$\Gamma$) is inductively defined as follows:
  \begin{enumerate}
  \item The type of $x$ is $!A$ iff $x:!A \in \Gamma$.
  \item If $N$ has type~$!B$ relative to $x:!A, \Gamma$, then the type of $\lambd[\typeof{x}{!A}][N]$ is $!A \lif !B$.
  \item If $N$ has type $!A \lif !B$ and $M$ has type~$!A$ then $(NM)$
  has type~$!B$
  \item If $N$ has type $!A$ and $M$ has type~$!B$, then $\pair{N}{M}$
  has type $!A \land !B$.
  \item If $N$ has type~$!A_1 \land !A_2$, then $\proj{i}{N}$ has type~$!A_i$.
  \item If $N$ has type $!A$, then
    $\inj[!B]{i}{!M}$ has type $!A \lor !B$ if $i=1$ and type $!B \lor
    !A$ if $i=2$.
  \item If $M$ has type $!A \lor !B$, $N_1$ has type $!C$ relative
  to $x_1:!A, \Gamma$, and $N_2$ has type~$!C$ relative to $x_2:!B,
  \Gamma$, then $\dcase{M}{x_1}{N_1}{x_2}{N_2}$ has type~$!C$.
  \item If $N$ has type~$\lfalse$, then
    $\abort{!A}{N}$ has type~$!A$.
  \end{enumerate}
\end{defn}

\begin{prop}
Every proof term~$N$ has exactly one type relative to a
context~$\Gamma$ for~$N$, if it has a type at all (i.e., is a correct
proof term).
\end{prop}

\begin{proof}
  By induction on~$N$.
\end{proof}

We write $\Gamma \Sequent N : !A$ if the type of $N$ relative to
$\Gamma$ is~$!A$. The clauses in \olref{defn:type} should seem familiar:
they are precisely the rules of \Log{tN2i} with one slight difference:
the context~$\Gamma$ is the same throughout. We can codify the
definition as a calculus \Log{tN3i} with the rules given in \olref{tab:tN2ip}.

\subfile{rules-tN3}

Proof terms are terms in a version of the $\lambd$-calculus, the
\emph{typed $\lambd$-calculus with products, sums, and empty type}. The
classical $\lambd$-calculus only has term built from variables,
\emph{$\lambd$-abstracts}~$\lambd[x][N]$ and application $(MN)$, and
it lacks the annotation by~$!A$ of the typed $\lambd$-abstract
$\lambd[\typeof{x}{!A}][N]$. The $\lambd$-calculus is a model of
computation which is Turing complete (i.e., every partial computable
function can be represented in it by a term). The version of the
$\lambd$-calculus given by proof terms and typing rules is a
restricted model of computation, but the basic ideas are the same. It
is restricted by the type assignment: only \emph{well-typed} terms,
i.e., proof terms with types (relative to a context~$\Gamma$) are
legal terms. For instance, $\lambd[x][(xx)]$ is a legal term in the
classical untyped $\lambd$-calculus, but
$\lambd[\typeof{x}{!A}][(xx)]$ is not well-typed for any $!A$. That's
because the typing rules require that in order for $(xx)$ to have a
type, the first~$x$ would have to have type $!A \lif !B$ and the
second $x$ type~$!A$, and that's impossible since $!A$ is a proper
sub-!!{formula} of~$!A \lif !B$.

We can think of types intuitively as the kind of thing named by a term
of that type. So a term of type $!A \land !B$ picks out a pair with
first component of type~$!A$ and second component of type~$!B$, i.e.,
!!a{element} of the product~$!A \times !B$. A term of type $!A \lif !B$
is a function from things of type $!A$ to things of type~$!B$. A term
of type $!A \lor !B$ is either something of type $!A$ or something of
type~$!B$, plus the information which one it is. This is like the
disjoint union of the two sets $!A$ and $!B$ defined as $\{\tuple{i,x}
: i = 1 \text{ or } 2, x \in !A \text{ if } i = 1, x \in !B \text{ if
} i=2\}$. The type $\lfalse$ is the empty set, i.e., if a term has
that type it can't name anything. Since natural deduction (including
\Log{tN3i}) is consistent, there is no !!{proof} of~$\lfalse$ unless
there are open assumptions, and so there is no closed proof term (term
without free variables) that has type~$\lfalse$.

The operators of the typed $\lambd$-calculus produce terms that
intuitively name the right sort of thing, provided the terms it is
applied to name certain things. E.g., ``if $N_1 : !A$ and $N_2 : !B$,
then $\pair{N_1}{N_2} : !A \land !B$'' says that if $N_1$ names
somethign of type~$!A$ and $N_2$ names something of type~$!B$, then
$\pair{N_1}{N_2}$ names something of type $!A \land !B$.
